\section{Discussion}
\label{sec:discussion}

\paragraph{Nominate the aggregate before the intervention.} Our result is only surprising if the
channel that motivated the fix is also the channel used to score it. It was, in our first pass, and
the conclusion we nearly reported --- ``look-ahead reduces the reward hack roughly fourfold'' --- is
true of over-praise, replicated by a second judge, mechanistically explained at three levels, and
wrong about the therapist's behaviour. The discipline that catches it is cheap: fix, in advance,
the aggregate the channel is a component of, and report the intervention against that. Here the
aggregate was already defined by the instrument (six coded behaviours summing to a session total),
which is a good argument for using coded behaviour inventories rather than bespoke channel metrics
as the outcome for a mitigation study.

\paragraph{Channel-level mitigations and the shape of the reward.} A reward that is hackable
through one route is usually hackable through several, because what makes it hackable is a gap
between what it measures and what is wanted, not the particular behaviour that exploits the gap.
Our rubrics measure how the session felt to the client. Flattery moves that; so, on this client
simulator, does confident advice. Closing the first route leaves the gap and the optimiser
untouched. On that reading, interventions that change \emph{what the grader sees} --- more context,
longer trajectories, richer transcripts --- should be expected to redistribute pressure, while
interventions that change \emph{what the grader is asked} --- adding an explicit MI-adherence
penalty to the training reward, rather than only to the evaluation --- are the ones with a route to
reducing the total. We did not test the latter and cannot claim it works.

\paragraph{Changing composition breaks aggregate comparability.} This is the part we did not
anticipate. Our two graders agree on the sign and the significance of every component of the
substitution, and disagree on the total, because they weight the substituted behaviour
differently: the held-out judge charges the $K{=}5$ policy $1.95$ acts of unsolicited advice where
the primary oracle charges $0.31$. When an intervention shifts \emph{which} failure occurs, the
aggregate stops being a grader-invariant quantity even though its parts remain ones. The
uncomfortable implication is that ``score the mitigation on the aggregate'' --- our own first
lesson --- is necessary but not sufficient: an aggregate over heterogeneous failures inherits
whatever relative weighting the grader happens to hold, and that weighting is invisible until
something moves the mix. Reporting the component vector alongside the total is the cheap fix, and
it is what we do in Table~\ref{tab:substitution}.

\paragraph{Trajectory rewards are not free.} Look-ahead cost more candidates and five extra
simulated turns each to reach the same iteration, and returned no improvement on the reward and no
grader-robust reduction in MI-inconsistency. Read strictly, on this task the correct summary is
that $K{=}5$ was not worth its price. Read more usefully, it did something specific, large and
verifiable to the reward's preference structure, and that capability may be worth having where the
target behaviour \emph{is} the aggregate of interest rather than one component of it.

\paragraph{The clinical reading.} From an MI standpoint the substitution is not neutral, even
though the count is. Over-praise and unsolicited advice are both MI-inconsistent, but they fail
differently: praise inflates the alliance while doing no therapeutic work, whereas advising without
permission asserts the therapist's agenda over the client's, which is the failure mode MI was
formulated against. Our judges disagree on whether the \Kf{} mix is more severe overall
(\S\ref{sec:aggregate}), so we do not rank them --- but we note that ``same number of violations''
should not be read as ``same clinical picture'', and that a count-based aggregate is itself an
approximation a clinician would want to interrogate.

\paragraph{What would change our mind.} A version of this result that reduced the total would have
to show the aggregate moving with the components, not instead of them. We would also update on a
second intervention with the same structure --- a different way of changing what the grader sees ---
producing a genuine aggregate reduction, which would make our result about look-ahead specifically
rather than about channel-level mitigation generally. We report the narrow claim.
